\section*{Chapitre \ref{ch:iv2:research-and-strategy-cases} --- Recherche et cas de stratégie}
\markboth{\protect\omTocMark{}{Réponses~: Recherche et cas de stratégie}}{\protect\omTocMark{}{Réponses~: Recherche et cas de stratégie}}

\begin{solution}{iq:iv2:research-and-strategy-cases:1}
C'est environ 1.06.

$0.0008/0.012 = 0.0667$ par jour, et $\times\sqrt{252} = 1.06$.

\iqwrong{Annualiser l'écart type par 252. La moyenne évolue avec le temps, et la volatilité avec sa racine carrée.}
\end{solution}

\begin{solution}{iq:iv2:research-and-strategy-cases:2}
La statistique $t$ vaut 2, et un ratio de Sharpe de 0.5 demande 16 ans.

Le $t$ est égal au ratio de Sharpe multiplié par la racine carrée du nombre d'années~: $1 \times \sqrt{4} = 2$. Puis $(2/0.5)^2 = 16$.

\iqwrong{Croire que des données quotidiennes apportent de la précision sur la moyenne. Qu'elles soient quotidiennes ou mensuelles, la longueur de l'échantillon fixe le $t$.}
\end{solution}

\begin{solution}{iq:iv2:research-and-strategy-cases:3}
Le ratio de Sharpe brut vaut environ 2.65, et le net environ 2.12.

Le coût vaut $0.4 \times 5 = 2$ points de base par jour. Donc $10/60 \times \sqrt{252} = 2.65$, et $8/60 \times \sqrt{252} = 2.12$.

\iqwrong{Comparer directement un coût à un ratio de Sharpe. Convertir d'abord le coût dans les mêmes unités quotidiennes.}
\end{solution}

\begin{solution}{iq:iv2:research-and-strategy-cases:4}
Cela fait 12.5~millions de dollars par titre, et 2.5~milliards de dollars sur les 200.

5~\% de 50~millions de dollars font 2.5~millions de dollars par jour, sur 5 jours.

\iqwrong{Supposer que la sortie est sans coût. La même limite s'applique à la sortie, et en situation de stress le volume obtenu n'est pas celui qu'on avait supposé.}
\end{solution}

\begin{solution}{iq:iv2:research-and-strategy-cases:5}
L'espérance vaut 0.1 par transaction, le ratio de Sharpe par transaction vaut environ 0.1, et sur 100 transactions il vaut environ 1.0.

La moyenne vaut $0.55 - 0.45 = 0.1$. L'écart type vaut $\sqrt{1 - 0.01} = 0.995$. Donc $0.1/0.995 \times 10 = 1.005$.

\iqwrong{Considérer un taux de réussite (hit rate) de 55~\% comme élevé. Un taux n'a pas de sens sans les tailles~: 55~\% de réussite avec des gains faibles et une perte importante fait perdre de l'argent.}
\end{solution}

\begin{solution}{iq:iv2:research-and-strategy-cases:6}
Il affiche un ratio de Sharpe annualisé d'environ 21. Avec un décalage d'un jour, il vaut zéro.

La position est le signe du rendement du jour, donc le P\&L vaut $|r|$. Sa moyenne vaut $\sqrt{2/\pi} = 0.798$, et son écart type $\sqrt{1 - 2/\pi} = 0.603$. Donc le ratio de Sharpe quotidien vaut 1.32, et $\times 15.87 = 21$.

Avec un décalage, on garde le signe d'aujourd'hui face au rendement indépendant de demain, et la moyenne vaut 0.

\iqwrong{Faire confiance à un ratio de Sharpe aussi élevé. Un résultat de 21 est un bug. Chercher une jointure de dates qui a utilisé des données du jour même.}
\end{solution}

\begin{solution}{iq:iv2:research-and-strategy-cases:7}
C'est environ 64~\%, et on en attend 1.

$1 - 0.95^{20} = 0.642$, et $20 \times 0.05 = 1$.

\iqwrong{Rapporter le meilleur test à son propre seuil de 5~\%. Le bon seuil pour le meilleur parmi 20 est d'environ 0.25~\%.}
\end{solution}

\begin{solution}{iq:iv2:research-and-strategy-cases:8}
C'est environ 16~\% pour un ratio de Sharpe de 1, et environ 2.3~\% pour un ratio de Sharpe de 2.

La moyenne divisée par l'écart type est le ratio de Sharpe, donc une perte a pour probabilité $\Phi(-\mathrm{SR})$.

\iqwrong{Attendre d'un ratio de Sharpe de 1 qu'il rapporte de l'argent chaque année. Une année sur six devrait être perdante, et un gérant de portefeuille doit survivre à cela.}
\end{solution}

\begin{solution}{iq:iv2:research-and-strategy-cases:9}
Il affiche 10~\%, alors que la réalité est d'environ 6.7~\%.

$0.97 \times 10~\% + 0.03 \times (-100~\%) = 6.7~\%$.

L'écart est donc de 3.3 points par an pour cette seule erreur, ce qui suffit à créer une stratégie à partir de rien. Les stratégies long small cap et distressed sont celles dont les résultats sont le plus gonflés.

\iqwrong{Utiliser les membres actuels de l'indice ou les titres cotés aujourd'hui comme univers du backtest. Utiliser la composition de l'indice à chaque date passée.}
\end{solution}

\begin{solution}{iq:iv2:research-and-strategy-cases:10}
On utilise un effet de levier de 0.4. Le ratio de Sharpe ne change pas, et le rendement baisse dans le même rapport que la volatilité.

$10/25 = 0.4$, en négligeant le financement de la trésorerie.

\iqwrong{Croire que l'effet de levier modifie le ratio de Sharpe. Il ne le modifie pas, avant financement et coûts.}
\end{solution}

\begin{solution}{iq:iv2:research-and-strategy-cases:11}
On attend un ratio de Sharpe d'environ 1.12. La règle approchée donne 1.36. Un ratio de Sharpe affiché de 1 est donc inférieur à ce que le pur bruit donne après 100 essais.

L'espérance du maximum de 100 normales centrées réduites vaut environ 2.51. Un ratio de Sharpe mesuré sur $T$ années a pour erreur type $1/\sqrt{T}$. Donc $2.51/\sqrt{5} = 1.12$.

La règle donne une borne supérieure, atteinte lentement à mesure que $N$ augmente~: $\sqrt{2\ln 100} = 3.03$, et $3.03/\sqrt{5} = 1.36$.

Il faut donc rapporter le nombre d'essais. C'est pourquoi un bon journal de recherche compte chaque variante, et pas seulement la version finale.

\iqwrong{Considérer les données hors échantillon (out-of-sample) comme à l'abri. Si on les a regardées plusieurs fois, elles sont aussi dans l'échantillon (in-sample).}
\end{solution}

\begin{solution}{iq:iv2:research-and-strategy-cases:12}
C'est environ 21 points de base par jour.

Avec une corrélation $\rho$ et des lois normales conjointes, la moyenne d'un décile du signal vaut $\rho\sigma$ fois la moyenne du signal dans ce décile. La moyenne du décile supérieur vaut $\varphi(1.2816)/0.1 = 1.755$. Donc le spread vaut $2 \times 1.755 \times 0.03 \times 2~\% = 0.21~\%$.

Cela vaut avant coûts et sans contrainte. Un IC de 0.03 est correct. Il paraît petit, mais le spread montre pourquoi il a sa valeur.

\iqwrong{Multiplier l'IC par la volatilité et s'arrêter là. Cela donne 6 points de base par écart type du signal. Le spread entre déciles doit atteindre environ 3.5 écarts types.}
\end{solution}

\begin{solution}{iq:iv2:research-and-strategy-cases:13}
La demi-vie vaut environ 6.6 jours. Le total est de 40 points de base. On capture environ 50~\% en 6.6 jours, et environ 88~\% en 20 jours.

$\ln 0.5/\ln 0.9 = 6.58$. Le total vaut $4/(1 - 0.9) = 40$. La part captée vaut $1 - 0.9^h$.

Donc la demi-vie indique un horizon de détention, et le total indique le coût qu'on peut supporter. Si un aller-retour coûte 10~points de base, trader ce signal en vaut largement la peine, mais toutes les petites variations du signal ne le justifient pas.

\iqwrong{Lire le coefficient quotidien comme la valeur du signal. Le chiffre quotidien paraît petit, mais c'est la somme qui rémunère la transaction.}
\end{solution}

\begin{solution}{iq:iv2:research-and-strategy-cases:14}
Avec un signal lent, la $t$-statistique usuelle est trop grande d'environ 4.4 fois, proche de $\sqrt{20} = 4.47$, de sorte qu'une hypothèse nulle vraie est rejetée environ 65~\% du temps au seuil de 5~\%. Si le signal change complètement chaque jour, la $t$-statistique est à peu près correcte. Il faut utiliser des erreurs de Newey-West avec au moins 20 retards, ou des données non chevauchantes.

Les rendements futurs partagent 19 observations sur 20 avec l'observation suivante. Lorsque le signal est aussi lent, le produit du signal et de l'erreur reste corrélé sur ces jours, et c'est ce produit qui fixe l'erreur type. Dans une simulation avec un coefficient AR de 0.99, la $t$ a un écart type de 4.4.

Avec un signal indépendant chaque jour, le produit n'est pas corrélé malgré le chevauchement, et l'erreur classique tient. Le problème exige donc à la fois le chevauchement et la persistance.

\iqwrong{Dire que le seul chevauchement gonfle la $t$-statistique. Le résidu est autocorrélé dans tous les cas, mais c'est son produit avec le régresseur qui compte.}
\end{solution}

\begin{solution}{iq:iv2:research-and-strategy-cases:15}
Cela fait environ 45 points de base à 5~\%, et environ 32 à 2.5~\%. Diviser l'ordre par deux réduit le coût total en dollars d'environ 65~\%. Le répartir sur deux jours le réduit de 29~\% au total.

$200\sqrt{0.05} = 44.7$, et $200\sqrt{0.025} = 31.6$. Donc par dollar, le coût baisse d'un facteur $\sqrt{2}$.

Le coût total est la taille multipliée par l'impact, il croît donc comme $q^{1.5}$. Une taille moitié moindre coûte $0.5^{1.5} = 35~\%$ du total. Trader la moitié sur chacun de deux jours coûte $2 \times 35~\% = 71~\%$.

C'est l'argument en faveur du fractionnement, tant que le signal reste valable. Il a son propre prix~: les coûts de délai, et un marché qui a vu une partie de l'ordre.

\iqwrong{Utiliser un modèle linéaire du coût, qui dit que la taille n'importe pas. Toute question de taille exige un impact concave par dollar et un total convexe.}
\end{solution}

\begin{solution}{iq:iv2:research-and-strategy-cases:16}
Il est maximal à 400 millions de dollars. L'avantage net y vaut 10 points de base, et le profit vaut 400\,000 dollars par unité de rotation. Le profit est nul à 900 millions de dollars.

Le profit par rotation vaut $Q(30 - 10\sqrt{Q})$. En annulant la dérivée~: $30 - 15\sqrt{Q} = 0$, donc $Q = 4$. Le coût vaut 20, et le net vaut 10 points de base sur 400 millions de dollars.

En général, à la meilleure taille, le coût absorbe les deux tiers de l'avantage brut. Et pour $Q = 9$, le coût égale l'avantage.

La capacité n'est donc pas un chiffre unique. Il y a la taille qui maximise les dollars, une taille plus petite qui maximise le ratio de Sharpe, et une taille plus grande à laquelle on ne gagne rien.

\iqwrong{Donner la taille à profit nul comme capacité. C'est là que l'on a cédé tout l'avantage au marché, et elle dépasse plus du double la taille raisonnable.}
\end{solution}

\begin{solution}{iq:iv2:research-and-strategy-cases:17}
Il faut 1, 4 et 16 ans. On ne prouvera donc pas en conditions réelles un ratio de Sharpe modeste, et la montée en charge doit reposer sur les limites de risque et sur l'adéquation des chiffres en production au modèle.

$t = \mathrm{SR}\sqrt{T}$, donc $T = (2/\mathrm{SR})^2$.

\textbf{Un calendrier.}
\begin{enumerate}
\item Papier ou petite taille, pendant quelques semaines~: vérifier que les exécutions, les coûts et les données correspondent au backtest.
\item 10~\% de la cible~: au moins 2 à 3 mois. Comparer le coût par dollar, la décroissance du signal et les taux d'exécution au modèle. Ces grandeurs se testent beaucoup plus vite que le P\&L~: elles ont bien plus d'observations.
\item 25~\% à 50~\%~: monter en charge lorsque les coûts et les expositions correspondent, et non parce que le P\&L est bon.
\item Cible, avec une limite de drawdown calibrée sur le ratio de Sharpe auquel on croit réellement, et une règle pour réduire la taille de moitié.
\end{enumerate}

\iqwrong{Attendre que le P\&L en conditions réelles le prouve, ou augmenter la taille après un bon mois. Les deux accordent trop d'importance au P\&L. Ce que l'on peut tester rapidement, c'est si le dispositif fonctionne comme prévu par le modèle.}
\end{solution}

\begin{solution}{iq:iv2:research-and-strategy-cases:18}
Il reste un point de base~: $20 - 11 - 3 - 2 - 2 - 1 = 1$. Le coût de délai et l'impact s'aggravent avec la taille, et le spread reste, par dollar, de même que le taux d'exécution manqué, sauf si l'on atteint une part importante du volume.

Pour le trouver, il faut vérifier les frais, l'emprunt et le financement, les prix de marquage, et le signal en production par rapport au même jour dans le backtest. Relancer le backtest sur les dates de production avec les ordres réels, et comparer ordre par ordre.

Les exécutions manquées comptent en général davantage qu'il n'y paraît~: on est exécuté quand le prix va contre la position, et on rate l'exécution quand il va dans son sens.

\iqwrong{Expliquer l'écart par un «~alpha decay~» sans détail. La décroissance vient en dernier, une fois que chaque coût mécanique a été compté.}
\end{solution}

\begin{solution}{iq:iv2:research-and-strategy-cases:19}
Le levier de Kelly vaut 15 fois, et la croissance logarithmique espérée est d'environ 113~\%. Si le rendement véritable est moitié moindre, on ne gagne rien en termes logarithmiques à ce levier. Au demi-Kelly, on gagne environ 28~\%.

$f = \mu/\sigma^2 = 0.15/0.01 = 15$. La croissance vaut $f\mu - \tfrac12 f^2\sigma^2 = 2.25 - 1.125 = 1.125$. Avec $\mu$ divisé par deux, $15 \times 0.075 - 1.125 = 0$. Au demi-Kelly, $7.5 \times 0.075 - 0.5 \times 56.25 \times 0.01 = 0.281$.

Donc le Kelly plein, appliqué à une estimation, rend exactement nul le gain d'une vraie stratégie. Comme la moyenne est la grandeur la plus difficile à estimer, n'utiliser qu'une fraction de Kelly n'a rien de timide~: c'est ce qu'il faut faire.

\iqwrong{Traiter Kelly comme une cible. C'est un plafond sur un avantage estimé, et il suppose l'absence de sauts et un rééquilibrage continu sans coût.}
\end{solution}

\begin{solution}{iq:iv2:research-and-strategy-cases:20}
Il faut d'abord vérifier que les données sont bien ce qu'elles prétendent être~: les prix, puis les opérations sur titres, puis le biais de survie, puis les horodatages. Le faire avant tout signal, et le montrer sur la première page.

\textbf{L'ordre des vérifications.}
\begin{enumerate}
\item \textbf{Forme.} Compter les actions par date, et les jours manquants. Repérer les jours fériés de marché et les trous.
\item \textbf{Cours aberrants.} Les plus fortes variations quotidiennes, avec une date et un ticker. S'inversent-elles le lendemain~? Vérifier si elles correspondent à un événement réel.
\item \textbf{Opérations sur titres.} Un split apparaît comme un rendement de $-50~\%$. Les prix sont-ils ajustés~? Qu'en est-il des dividendes~?
\item \textbf{Biais de survie.} Des titres sortent-ils du fichier~? Si les 500 sont présents pendant les cinq ans, l'échantillon a été choisi avec le recul, et il faut dire quel biais cela introduit.
\item \textbf{Horodatages.} À quel moment une valeur devient-elle connue~? Un cours de clôture est-il utilisable ce jour-là~?
\item \textbf{Cohérence.} Comparer la moyenne transversale à un indice connu, si l'on en a un. Distribution des rendements et corrélations.
\end{enumerate}

\textbf{La rédaction.} Un tableau de ce qui a été trouvé, modifié et laissé en l'état. Présenter un avant et un après pour une erreur claire. Les relecteurs retiennent cela bien plus qu'un modèle astucieux.

\textbf{Comment construire le sien.} Tenir une liste de contrôle sous forme de code que l'on peut relancer. Un second jeu de données demande les mêmes contrôles.

\iqwrong{Commencer par un modèle. Un ratio de Sharpe élevé trouvé à partir d'une erreur est l'échec classique d'un travail à domicile.}
\iqfollow{Vous trouvez 12 cours présentant des variations supérieures à 50~\%. Les supprimez-vous~?}{On ne les supprime pas sans les examiner. Les événements réels restent, car ils correspondent au rendement qu'a réellement obtenu un détenteur. Il faut corriger ou signaler ce qui est faux, et montrer le résultat dans les deux cas.}
\end{solution}

\begin{solution}{iq:iv2:research-and-strategy-cases:21}
Le meilleur mélange est de 80~\% et 20~\%, avec un IC d'environ 0.0306. À une corrélation de 0.667, le second signal ne reçoit aucun poids. Un signal faible peu corrélé peut valoir davantage qu'un signal fort mais redondant.

On résout $Cw = \mathrm{IC}$~: $w \propto (0.0267, 0.0067)$. L'IC combiné vaut $\sqrt{\mathrm{IC}^\top C^{-1}\mathrm{IC}} = 0.0306$.

Le poids du second est nul lorsque $\rho = 0.02/0.03 = 0.667$.

À une corrélation nulle, la valeur combinée vaudrait $\sqrt{0.03^2 + 0.02^2} = 0.036$. C'est nettement plus que 0.0306. Il faut donc vérifier l'IC résiduel d'un nouveau signal après les signaux existants, et non son IC isolé.

\iqwrong{Ajouter des signaux à poids égal, ou additionner les IC. Les IC ne s'additionnent pas, et un signal corrélé apporte très peu.}
\end{solution}

\begin{solution}{iq:iv2:research-and-strategy-cases:22}
Le P\&L est de 20~millions de dollars, et net des coûts, de 10~millions de dollars, soit 5~\% du capital. Le ratio de Sharpe de point mort est de 0.5.

10~\% de 200~millions de dollars font 20~millions de dollars de risque, et un ratio de Sharpe de 1 rapporte 20~millions de dollars.

\textbf{Ce qui me ferait changer d'avis.} Ce ratio de Sharpe est-il net des coûts de transaction~? Le capital de la société est-il en drawdown~? Les coûts peuvent-ils baisser, ou augmentent-ils avec la taille~? Un rendement de 5~\% est-il supérieur à ce que coûte le capital de la société~? Avec un taux de hurdle de 8~\%, il faudrait 16~millions de dollars de plus que le point mort, soit un ratio de Sharpe de 1.3.

\iqwrong{Comparer 10~millions de dollars à zéro. Le capital a un coût, et un desk qui rapporte moins que ce coût réduit la taille de la société.}
\end{solution}

\begin{solution}{iq:iv2:research-and-strategy-cases:23}
Le chiffre d'affaires est d'environ 730\,000~\$ par an, ce qui ne couvre pas 2~millions de dollars. Le point mort correspond à une part d'environ 5.5~\%. Le plan tel qu'il est énoncé ne fonctionne donc pas, à moins que la part ou la captation nette n'augmente fortement.

$10^9 \times 0.02 \times 10^{-4} \times 365 = \$730\,000$. D'où $2\,000\,000/(10^9 \times 10^{-4} \times 365) = 5.48~\%$.

Cela indique le chiffre à remettre en cause. Est-ce 1~point de base net de sélection adverse~? Un 5~\% est-il réaliste face à ceux qui sont déjà présents~? Peut-on ajouter des marchés en réutilisant la même infrastructure~?

\iqwrong{Multiplier par un gros volume total et s'en étonner. La part et la captation nette sont les deux chiffres difficiles, et la plupart des projets supposent les deux.}
\end{solution}

\begin{solution}{iq:iv2:research-and-strategy-cases:24}
L'espérance est de 0.4 par transaction, l'écart type d'environ 1.91, et le ratio de Sharpe sur 100 transactions d'environ 2.1. Une série de 10 pertes se produit dans environ 1.3~\% des cas.

$0.35 \times 3 - 0.65 = 0.4$. La variance vaut $0.35 \times 9 + 0.65 - 0.16 = 3.64$. Donc $0.4/1.91 = 0.21$ par transaction, et $\times 10 = 2.1$.

$0.65^{10} = 1.35~\%$. Sur une carrière comportant beaucoup de transactions, on le verra souvent.

\iqwrong{Supposer que des transactions indépendantes sont réalistes. Les positions de tendance détenues en même temps sont corrélées, si bien que le ratio de Sharpe réel est bien inférieur à 2.1.}
\end{solution}

\begin{solution}{iq:iv2:research-and-strategy-cases:25}
L'alpha vaut 11~\%, la volatilité résiduelle 6~\%, et le ratio d'information environ 1.83.

$15 - 0.5 \times 8 = 11$. La variance résiduelle vaut $0.01 - 0.25 \times 0.0256 = 0.0036$.

Le ratio de Sharpe de 1.5 mélange l'exposition au marché et le talent. Un desk peut acheter du bêta pour presque rien, et c'est le résiduel qu'il paie. Ici, le ratio d'information est supérieur au ratio de Sharpe, car la partie bêta n'a qu'un ratio de Sharpe de 0.5.

\iqwrong{Comparer des stratégies par ratio de Sharpe alors que certaines portent une exposition au marché. Dans un marché haussier, le bêta ressemble à du talent.}
\end{solution}

\begin{solution}{iq:iv2:research-and-strategy-cases:26}
Il faut d'abord demander combien de variantes ont été testées. Il faut ensuite s'interroger sur la période, le modèle de coûts, l'hypothèse d'exécution et l'univers. Un ratio de Sharpe de 3.5 sur deux ans doit inciter à chercher ce qui ne va pas avant toute autre chose.

\textbf{Dans l'ordre.}
\begin{enumerate}
\item \textbf{Essais.} Combien de spécifications~? Avec 100 spécifications, un ratio de Sharpe dans l'échantillon de 1.8 sur 2 ans provient du bruit.
\item \textbf{Période.} Quelles deux années~? Comprend-elle un épisode de stress~? Quel est le ratio de Sharpe dans chaque moitié, et par mois~?
\item \textbf{Exécution à la clôture.} Le signal est-il calculé sur cette même clôture~? Si oui, il y a un biais d'anticipation. L'enchère de clôture est-elle suffisamment importante~?
\item \textbf{Coûts.} Les 5~points de base sont-ils réalistes pour ce taux de rotation et pour cette taille~? Quel est le coût de point mort~?
\item \textbf{Univers.} Grandes capitalisations d'aujourd'hui, ou à chaque date~?
\item \textbf{Expositions.} Bêta, secteur, facteurs. Que reste-t-il après les avoir retirés~?
\end{enumerate}

Chaque réponse restreint le champ. Si la stratégie résiste à tout cela, je voudrais la déployer en production à petite taille, car les deux années suivantes sont le seul test propre.

\textbf{Comment construire sa propre version.} Il faut conserver cette liste dans un ordre fixe, afin de poser les mêmes questions à son propre travail.

\iqwrong{Attaquer d'abord le modèle. Les problèmes habituels se trouvent dans les données et dans le processus, bien avant la formule.}
\end{solution}

\begin{solution}{iq:iv2:research-and-strategy-cases:27}
Son ratio de Sharpe sur la première moitié vaut environ 1.59, sur 2.5 ans, et son ratio de Sharpe sur la seconde moitié vaut environ 0, avec un rang proche de 50 sur 100. Le rang dans l'échantillon ne vous apprend donc rien.

Le maximum vaut 2.51 erreurs types, et l'erreur type pour $T = 2.5$ vaut $1/\sqrt{2.5} = 0.632$. D'où 1.59.

Sans talent, les deux moitiés sont indépendantes, si bien que la stratégie retenue n'est qu'un tirage aléatoire sur l'autre moitié. Son ratio de Sharpe espéré vaut 0, et son rang est uniforme~: il se situe sous la médiane une fois sur deux.

\textbf{Le test.} On découpe les données de nombreuses façons. À chaque fois, on retient la meilleure stratégie dans l'échantillon et on note son rang hors échantillon. La probabilité de surapprentissage du backtest est la part des découpages où ce rang se situe sous la médiane. Une valeur d'environ 50~\% signifie que vous avez sélectionné du bruit. Une valeur faible signifie que le rang dans l'échantillon porte de l'information.

\iqwrong{Utiliser un seul découpage, car un seul découpage ne donne qu'un seul tirage. Et ne pas réutiliser la moitié hors échantillon pour choisir, sinon elle devient de l'échantillon.}
\end{solution}

\begin{solution}{iq:iv2:research-and-strategy-cases:28}
Il vaut environ 4~\% dans l'échantillon, et environ $-4~\%$ hors échantillon. Un $R^2$ dans l'échantillon de 4~\% avec 20 variables et 500 observations est donc ce que produit le pur bruit.

Dans l'échantillon, les moindres carrés ajustent environ $k/T = 20/500$ du bruit. Hors échantillon, les coefficients ajustés sont eux-mêmes du bruit, et ajoutent de la variance à la prévision~: un autre $k/T$. L'erreur quadratique moyenne dépasse donc celle d'une prévision nulle, et le $R^2$ est négatif.

Une lecture honnête utilise le $R^2$ ajusté~: environ 0 ici.

\iqwrong{Penser que 4~\% est une affirmation modeste, donc petite. Sur les marchés, 4~\% serait énorme. Sa taille indique donc sans doute qu'il s'agit d'un surajustement.}
\end{solution}

\begin{solution}{iq:iv2:research-and-strategy-cases:29}
La moyenne vaut 12~\%, la volatilité environ 20.4~\%, et le ratio de Sharpe environ 0.59.

Moyenne en régime calme~: $2 \times 10~\% = 20~\%$. En régime de stress~: $-1 \times 20~\% = -20~\%$. Donc $0.8 \times 20 - 0.2 \times 20 = 12$.

La variance à l'intérieur des régimes vaut $0.8 \times 0.01 + 0.2 \times 0.04 = 0.016$. La variance entre les régimes vaut $0.8 \times 0.08^2 + 0.2 \times 0.32^2 = 0.0256$. La variance totale vaut 0.0416.

Le ratio de Sharpe est inférieur à la moyenne pondérée par le temps, de 1.4, car le déplacement de la moyenne entre les régimes ajoute de la variance. Un backtest sur le seul régime calme affiche 2, avec la moitié du risque.

Je le mesurerais donc par régime, je le réduirais lorsqu'une mesure de stress augmente, et je fixerais la limite à partir du chiffre du régime de stress.

\iqwrong{Moyenner les ratios de Sharpe entre régimes. Les ratios de Sharpe ne se moyennent pas. Il faut d'abord combiner les moyennes et les variances.}
\end{solution}

\begin{solution}{iq:iv2:research-and-strategy-cases:30}
Le profit est de 3.8~\% par transaction, la durée d'environ 21 jours, et si l'on était toujours en position, le rendement serait d'environ 47~\% par an. Cela ne tient pas compte de la rupture de la relation.

Le profit brut vaut $2 \times 2~\% = 4~\%$, moins 0.2~\%.

La réversion à la moyenne donne $\kappa = \ln 2/10 = 0.0693$. Un chemin moyen met un temps infini à atteindre exactement zéro, mais le bruit permet d'y arriver. Le résultat de premier passage admet une forme fermée~: la probabilité de ne pas encore avoir atteint zéro à partir de $x_0$ après le temps $t$ vaut $\mathrm{erf}\bigl(x_0 e^{-\kappa t}/(\sigma\sqrt{2(1 - e^{-2\kappa t})})\bigr)$, avec $\sigma$ l'écart type stationnaire. À partir de 2 écarts types, on obtient une durée moyenne d'environ 20.6 jours et une médiane d'environ 16.5. D'où environ 12.3 transactions par an. On obtient alors $3.8~\% \times 12.3 = 47~\%$.

\textbf{Non pris en compte.} Attendre 2 écarts types prend du temps, si bien que l'on n'est pas toujours en position. La demi-vie et la moyenne sont estimées, et peuvent se déplacer. Le risque le plus important survient lorsqu'un écart est large parce que la relation est rompue, par exemple lors d'une fusion, d'une faillite ou de l'arrivée d'un nouveau concurrent. Un modèle de réversion à la moyenne vous pousse alors à renforcer votre position au pire moment. Il faut donc fixer un stop, ou une limite de durée, à partir de la demi-vie.

\iqwrong{Utiliser le temps déterministe, $\ln(2\sigma/\epsilon)/\kappa$, qui diverge en zéro. Le bruit compte. Il faut simuler, ou utiliser le résultat de premier passage.}
\end{solution}

\begin{solution}{iq:iv2:research-and-strategy-cases:31}
La moyenne vaut environ 11~\%, la volatilité environ 12.3~\%, et le ratio de Sharpe environ 0.89, au lieu de 1.2. Il faut environ 114 ans pour l'observer avec une probabilité de 90~\%.

Moyenne~: $0.98 \times 12~\% + 0.02 \times (-40~\%) = 10.96~\%$. Variance~: $0.98(0.01 + 0.0144) + 0.02 \times 0.16 - 0.1096^2 = 0.0151$.

$1 - 0.98^n \ge 0.9$ signifie $n \ge \ln 0.1/\ln 0.98 = 114$.

Aucun backtest ne permet donc de l'écarter. Il faut raisonner à partir de la structure~: qu'est-ce qui ferait perdre 40~\% à cette stratégie~? Vend-elle de l'assurance~? Est-elle levée~? Il faut alors garder du capital en réserve contre ce scénario, grâce à un test de résistance plutôt qu'à une estimation.

\iqwrong{Affirmer que le backtest montre l'absence d'un tel risque. Un backtest ne peut pas montrer ce qui ne s'est pas produit. Les données les plus longues permettent seulement de dire que la fréquence de l'événement rare est faible, et non qu'elle est nulle.}
\end{solution}

\begin{solution}{iq:iv2:research-and-strategy-cases:32}
L'avantage net vaut environ 0.68 point de base à 50 millisecondes, 0 à 200 et $-0.5$ à 400. La latence de point mort est de 200 millisecondes. Réduire la latence de 200 à 50 millisecondes fait donc passer l'avantage de zéro à 0.68.

$2 \times 2^{-L/200}$~: 1.68, 1.00 et 0.50, moins 1 de coût.

\textbf{Ce que je vérifierais d'abord.} Où passent les 200 millisecondes~? Les concurrents sont-ils à 10 millisecondes, auquel cas le signal mesuré a déjà disparu au moment où vous le voyez~? La latence jusqu'à l'exécution compte, et pas seulement celle jusqu'à l'envoi. La décroissance du signal est une moyenne~: aux instants où il est le plus fort, il pourrait aussi se dissiper le plus vite.

Réduire la latence n'est un investissement rentable que si l'on conserve les exécutions. Si des rivaux plus rapides captent les meilleures transactions, on ne récupère que les restes.

\iqwrong{Supposer qu'une courbe de décroissance mesurée à partir de vos propres exécutions passées reflète le vrai signal. Ces exécutions sont déjà sélectionnées. Mesurez-la à partir du flux de cotations.}
\end{solution}

\begin{solution}{iq:iv2:research-and-strategy-cases:33}
La moyenne reste inchangée, proche de 0, la volatilité baisse légèrement, d'environ 4.47~\% à environ 4.28~\%, et le stop se déclenche environ 21~\% du temps. Un stop aide avec le momentum, et il nuit avec la réversion à la moyenne.

Sur une martingale, arrêter selon une règle quelconque conserve l'espérance~: c'est le théorème d'arrêt optionnel. Ainsi, sur une marche aléatoire, le stop ne peut que remodeler la distribution~: il coupe certaines pertes en queue de distribution, au prix de manquer certains rebonds.

Ce n'est donc pas une source de rendement. C'est un outil de gestion du risque lorsqu'on dispose d'une limite. Si les rendements suivent une tendance, une perte précoce annonce des pertes supplémentaires, et le stop aide. S'ils reviennent à la moyenne, le stop vend au pire moment.

Par ailleurs, les sauts de cours peuvent franchir un stop, si bien que sa protection est limitée. Et les coûts augmentent.

\iqwrong{Tester des stops sur un backtest et se réjouir de la baisse du drawdown. Regarder le ratio de Sharpe, et se demander si la règle détecte une véritable tendance ou se contente de réduire le risque.}
\end{solution}

\begin{solution}{iq:iv2:research-and-strategy-cases:34}
La validation croisée aléatoire à cinq plis donne un $R^2$ d'environ $+15~\%$ sur du pur bruit. Les blocs donnent environ $-46~\%$, et les blocs purgés environ $-55~\%$. Seules les versions par blocs révèlent qu'il n'y a rien à trouver.

\textbf{Pourquoi.} La ligne de test de chaque pli aléatoire se trouve à côté de lignes d'apprentissage aux variables presque identiques, et les étiquettes partagent 19 jours sur 20. Le modèle mémorise le bruit voisin, et les lignes de test le récompensent.

\textbf{Les blocs.} Ils séparent les données de test dans le temps, si bien que les lignes ne sont pas voisines. Elles le restent néanmoins aux bords des blocs, et une purge de 20 jours supprime ces lignes. Le $R^2$ fortement négatif montre ce que dix variables persistantes peuvent ajuster sur une courte période, et à quel point cela se dégrade sur une nouvelle période.

Donc, en finance, il faut utiliser des plis par blocs et purgés, avec un embargo après chaque bloc de test.

\iqwrong{Faire confiance à un $R^2$ élevé en validation croisée sur des étiquettes qui se chevauchent. Les plis aléatoires conviennent aux lignes indépendantes. Les séries temporelles qui se chevauchent doivent séparer strictement le test de l'apprentissage.}
\end{solution}

\begin{solution}{iq:iv2:research-and-strategy-cases:35}
Il faut environ 2\,400 transactions pour un $t$ espéré de 1.96, et environ 4\,900 pour une puissance de 80~\%. À 51~\%, il en faut quatre fois plus~: environ 9\,600 et 19\,600.

L'avantage par transaction vaut $2p - 1 = 0.04$, avec un écart type proche de 1. On a donc $t = 0.04\sqrt{N}$, et $N = (1.96/0.04)^2 = 2\,401$.

Pour une puissance de 80~\%, il faut que la moyenne se situe à $1.96 + 0.84$ erreurs types de 0~: $(2.8/0.04)^2 = 4\,900$.

Diviser l'avantage par deux multiplie la taille d'échantillon par quatre. Le premier nombre correspond à une puissance de 50~\%~: une fois sur deux, on manquerait encore l'effet.

\iqwrong{Traiter de nombreuses transactions corrélées comme si elles étaient indépendantes. Les transactions du même jour sur le même marché comptent beaucoup moins.}
\end{solution}

\begin{solution}{iq:iv2:research-and-strategy-cases:36}
On obtient $a = 30.9~\%$ et $b = 1.09$ (avec $Q$ en millions), et une prédiction d'environ 6.5~\% pour 500 millions de dollars. La taille optimale est d'environ 250 millions de dollars si le capital coûte 5~\%.

Les deux équations sont $20 = a - 10b$ et $12 = a - 17.32b$. D'où $b = 8/7.32 = 1.093$, et $a = 30.93$.

Gain net en dollars~: $Q(a - b\sqrt{Q})$. Son rendement marginal vaut $a - 1.5b\sqrt{Q}$. En l'égalant à 5~: $\sqrt{Q} = 25.93/1.64 = 15.8$, donc $Q = 250$.

Le marginal est ce que rapporte le dollar suivant. À 300 millions de dollars, la moyenne est de 12~\%, mais le dernier dollar ne rapporte qu'environ $30.93 - 28.4 = 2.5~\%$, soit moins que le coût. Une moyenne qui reste satisfaisante peut masquer ce phénomène.

\iqwrong{Augmenter la taille tant que le rendement moyen reste au-dessus du seuil de rentabilité. On ajoute alors du capital au-delà du point où les dollars supplémentaires rapportent moins qu'ils ne coûtent.}
\end{solution}

\begin{solution}{iq:iv2:research-and-strategy-cases:37}
Elle vaut environ 3.5 millions de dollars par an si le vendeur a raison, et environ 0.9 million si le vrai ratio de Sharpe n'est que la moitié~: moins que le prix. Il faut donc la tester soigneusement avant de payer.

Le ratio de Sharpe combiné de stratégies non corrélées vaut $\sqrt{1 + 0.25} = 1.118$. Pour 30 millions de dollars de risque, le gain vaut $0.118 \times 30 = 3.54$ millions de dollars. Avec 0.25~: $\sqrt{1.0625} = 1.031$, soit $0.92$ million de dollars.

\textbf{Avant de payer.} Obtenir un essai sur des données mises de côté. Vérifier que les horodatages des données correspondent bien à la date concernée, et qu'ils ne sont pas révisés. Se renseigner sur la durée de commercialisation et sur le nombre de clients. Mesurer son IC après les signaux existants. Vérifier que la période d'échantillon du vendeur n'inclut pas celle qui a conduit à son lancement.

\iqwrong{Comparer le ratio de Sharpe isolé du vendeur au prix. Ce que l'on achète, c'est l'augmentation, qui croît comme le carré du nouveau ratio de Sharpe et dépend de la corrélation avec le portefeuille existant.}
\end{solution}

\begin{solution}{iq:iv2:research-and-strategy-cases:38}
Mon estimation donne un vrai ratio de Sharpe d'environ 0.35 après le backtest, et d'environ 0.36 après 6 mois en production, avec un écart type d'environ 0.37. La stratégie n'est pas cassée~: le chiffre en production s'accorde avec cette estimation. Le backtest était surévalué. Je la ferais tourner à une taille fixée d'après 0.36, et non 2.4, et je l'arrêterais sur un drawdown défini à l'avance ou sur un changement de ses mécanismes.

\textbf{Étape 1~: la décote.} Le meilleur de 200 essais nuls sur 3 ans a un ratio de Sharpe espéré de $2.746/\sqrt{3} = 1.585$. Il ne reste donc à expliquer que $2.4 - 1.585 = 0.815$.

\textbf{Étape 2~: la loi a priori.} Traitons 0.815 comme une observation bruitée, de variance $1/3$, sous une loi a priori de variance 0.25. La variance a posteriori vaut $1/(4 + 3) = 0.143$, et sa racine carrée 0.378. La moyenne vaut $0.143 \times 0.815 \times 3 = 0.349$.

\textbf{Étape 3~: la période en production.} Six mois correspondent à une variance de 2, avec une observation de 0.5. La nouvelle précision vaut $7 + 0.5 = 7.5$, donc la variance vaut 0.133, et sa racine carrée 0.365. La moyenne vaut $(7 \times 0.349 + 0.5 \times 0.5)/7.5 = 0.359$.

Six mois d'observations déplacent donc à peine cette estimation~: six mois ne disent pas grand-chose. La probabilité que le vrai ratio de Sharpe soit négatif est d'environ 16~\%.

\textbf{La stratégie est-elle cassée~?} Un 0.5 en production est très inférieur à 2.4 en valeur, mais six mois de données ont une erreur type proche de 1.4. Elle n'est donc qu'à environ 0.1 erreur type de 0.36, et à environ 1.3 de 2.4~: le chiffre en production ne permet pas de les départager. La bonne question n'est pas ce qu'a rapporté le P\&L. Il s'agit de vérifier que les coûts et les exécutions correspondent, et que le signal prédit au même rythme que dans le backtest.

\textbf{Taille.} Une taille conçue pour 2.4 serait presque sept fois trop grande. Pour un ratio de Sharpe de 0.36 à un budget de risque donné, le levier demi-Kelly est faible, et dans un portefeuille, c'est ce qu'elle ajoute à l'ensemble qui compte. Je commencerais avec une part de risque faible, disons 15~\% de ce que suggère le backtest.

\textbf{Règle d'arrêt.} Fixée à l'avance~: un drawdown d'environ 2 à 2.5 années de rendement espéré, un changement net dans les mécanismes, une défaillance du modèle de coûts, ou une chute de l'IC mesuré du signal vers zéro. Pas à cause d'un mauvais mois.

\textbf{Au gérant.} «~Le backtest à 2.4 n'est pas le bon chiffre. Après la recherche et avec une loi a priori sceptique, j'estime le ratio à environ 0.35, plus ou moins 0.4, si bien qu'il a 16~\% de chances de ne pas avoir d'avantage. Je la garderais petite, je passerais les six prochains mois sur les exécutions et les coûts, puis je réexaminerais. Voici la règle qui me ferait l'arrêter.~»

\iqwrong{Comparer le live au backtest et conclure qu'elle est cassée, ou qu'elle va bien. Aucune de ces conclusions n'est possible sur six mois. Ce qui a modifié votre avis, c'est le nombre d'essais, et c'est ce que la plupart des gens omettent.}
\iqfollow{Et si les 200 essais étaient des variantes fortement corrélées~?}{Alors le nombre effectif est plus petit, la décote baisse, et l'estimation monte. Il faut estimer le nombre effectif à partir des corrélations, et présenter les deux.}
\end{solution}

\begin{solution}{iq:iv2:research-and-strategy-cases:39}
Le chiffre d'affaires vaut environ 3.65 millions de dollars, 7.3 millions et 18.25 millions pour des parts de 1~\%, 2~\% et 5~\%. La part de point mort vaut environ 1.7~\% sans la charge en capital, et environ 2.5~\% avec elle~; avec une capture moitié moindre, elle vaut environ 3.4~\% sans la charge, et environ 5.0~\% avec. Je ne m'y engagerais que sur la preuve du taux de capture et sur une voie vers 3~\%, avec un plan clair d'arrêt.

\textbf{Chiffre d'affaires.} $\$20 \text{ milliards} \times s \times 0.5 \times 10^{-4} \times 365 = \$365 \text{ millions} \times s$.

\textbf{Coûts.} Des frais fixes de 6 millions de dollars, et un risque de plateforme de $0.9~\% \times \$20 \text{ millions} = \$180\,000$. Le capital à 15~\% vaut 3 millions de dollars.

\textbf{Point mort.} $6.18/365 = 1.69~\%$. Avec le capital, $9.18/365 = 2.52~\%$. Avec une capture moitié moindre, $6.18/182.5 = 3.39~\%$, et avec le capital, $9.18/182.5 = 5.03~\%$.

\textbf{Ce que je demanderais.}
\begin{itemize}
\item Captation~: les 0.5 point de base sont-ils nets de la sélection adverse, des frais, de la perte sur inventaire et du financement~? Sur quels marchés~? Comment a-t-on mesuré cette captation~?
\item Part de marché~: qui opère déjà sur ces marchés~? Une part de 2~\% suppose de la prendre à quelqu'un qui ne la cédera pas.
\item Rabais et paliers de frais~: ils font souvent la différence, et une plateforme peut les modifier.
\item Inventaire~: 20 millions de dollars répartis sur de nombreuses plateformes, c'est peu. Faut-il davantage de capital, et reste-t-il disponible en situation de stress~?
\item Équipe~: quel est le plan si deux personnes clés partent~?
\end{itemize}

\textbf{Ce qui me ferait dire non.}
\begin{itemize}
\item Si le taux de captation provient d'un backtest, et qu'aucune donnée en production n'existe.
\item Si le plan suppose à la fois une part supérieure à 3~\% et une captation au haut de la fourchette.
\item Si la défaillance d'une plateforme effacerait tout le capital~: il faut limiter l'exposition, ou ne pas démarrer.
\item Si une part importante du chiffre d'affaires repose sur le rabais d'une seule plateforme.
\end{itemize}

\textbf{Que faire ensuite.} Financer un pilote réduit, avec un budget et une échéance~: faire tourner la stratégie sur deux marchés pendant six mois, et mesurer la captation, la part et la sélection adverse. Ces trois chiffres décident de tout. Arrêter si la captation est inférieure à la moitié du plan.

\iqwrong{Construire un tableur volumineux, avec des prévisions de chiffre d'affaires détaillées sur cinq ans. Un recruteur veut les deux ou trois chiffres dont dépend le résultat, et ce qui ferait changer votre décision.}
\end{solution}

\begin{solution}{iq:iv2:research-and-strategy-cases:40}
Le rendement net est meilleur autour de 5 à 6 jours, avec environ 0.94 point de base par jour et un ratio de Sharpe d'environ 1.33. Un rafraîchissement quotidien perd de l'argent après coûts, et un rafraîchissement à 20 jours rapporte environ 0.57 point de base par jour, avec un ratio de Sharpe d'environ 0.8.

\textbf{Dérivation.} Le coefficient AR vaut $\rho = 2^{-1/5} = 0.8706$. Avec des positions $s_{t_0}/E|s|$ normalisées à un brut égal à 1, la normalisation par action fait intervenir $E|s| = \sqrt{2/\pi}$.

\textit{Alpha brut.} La position d'âge $k$ rapporte $2\rho^k/E|s|$ points de base. En moyennant sur $k = 0, \dots, h - 1$,
\[
\alpha(h) = 2\sqrt{\pi/2}\;\frac{1 - \rho^h}{h(1 - \rho)}.
\]

\textit{Rotation.} Le rafraîchissement déplace $s_{t} - s_{t-h}$, qui est gaussienne de variance $2(1 - \rho^h)$. Sa taille absolue moyenne, rapportée à $E|s|$, vaut donc $\sqrt{2(1 - \rho^h)}$, par période de $h$ jours.

\textit{Coût.} $5\sqrt{2(1 - \rho^h)}/h$ points de base par jour.

\textit{Volatilité.} Avec 500 actions indépendantes, $\sqrt{\pi/2} \times 200/\sqrt{500} = 11.2$ points de base par jour.

\begin{center}
\begin{tabular}{rrrrr}
\toprule
$h$ & Alpha (pb/jour) & Coût (pb/jour) & Net & Sharpe \\
\midrule
1 & 2.51 & 2.54 & $-0.04$ & $-0.05$ \\
3 & 2.20 & 1.37 & 0.82 & 1.16 \\
5 & 1.94 & 1.00 & 0.94 & 1.33 \\
6 & 1.82 & 0.89 & 0.94 & 1.33 \\
10 & 1.45 & 0.61 & 0.84 & 1.19 \\
20 & 0.91 & 0.34 & 0.57 & 0.80 \\
\bottomrule
\end{tabular}
\end{center}

Chaque colonne est arrondie séparément. Une simulation du portefeuille complet confirme le tableau.

\textbf{Ce que je rédigerais.}
\begin{itemize}
\item Le résultat principal~: le meilleur rafraîchissement se situe autour d'une semaine, avec un ratio de Sharpe net d'environ 1.3, et un ratio de Sharpe brut d'environ 2.7. Les coûts absorbent donc environ la moitié du signal.
\item La courbe est plate entre 3 et 10 jours. Je ne chercherais donc pas à affiner $h$~; je retiendrais une valeur intermédiaire.
\item Ce qui manque~: les coûts sont ici linéaires, sans impact~; le bruit est indépendant d'une action à l'autre, ce qui gonfle le ratio de Sharpe. Avec un facteur de marché, une bonne partie disparaîtrait.
\item Une meilleure règle qu'un rafraîchissement fixe consiste à se déplacer, chaque jour, en partie vers la cible, ce qui fait mieux que n'importe quel $h$ fixe lorsque le coût est linéaire. C'est ce que j'essaierais ensuite, et je le dirais.
\end{itemize}

\iqwrong{Présenter le ratio de Sharpe brut à $h = 1$, d'environ 3.5, comme la qualité du signal. Les coûts font partie de la valeur du signal~: avec un rafraîchissement quotidien, ce signal ne rapporte rien.}
\end{solution}
